\begin{corollary}[1.8.7]
\label{ega2:addendum:I.1.8.7-fr}
Bajo las hipótesis de \sref{ega2:addendum:I.1.8.6-fr}, para todo $A$-módulo $M$ y todo
$\sh{O}_X$-módulo $\sh{F}$ existe un isomorfismo canónico, funtorial
en $M$ y $\sh{F}$,
\begin{equation}
\label{ega2:addendum:I.1.8.7.1-fr}
  \Hom_{\sh{O}_S}(\widetilde M,u_*(\sh{F}))
  \xrightarrow{\sim}\Hom_A(M,\Gamma(X,\sh{F})).
\tag{1.8.7.1}
\end{equation}
\end{corollary}

\begin{proof}
Por (\textbf 0, 4.4.3) y el Lema~\sref{ega2:addendum:I.1.8.5-fr}, existe un isomorfismo
canónico de bifuntores
\[
  \Hom_{\sh{O}_S}(\widetilde M,u_*(\sh{F}))
  \xrightarrow{\sim}\Hom_{A'}(M',i_*(u_*(\sh{F})))
\]
El segundo miembro es claramente $\Hom_A(M,\Gamma(X,\sh{F}))$. Bajo
el homomorfismo canónico \sref{ega2:addendum:I.1.8.7.1-fr}, a un
$\sh{O}_S$-homomorfismo
$h:\widetilde M\to u_*(\sh{F})$ (equivalentemente, a un $u$-morfismo
$\widetilde M\to\sh{F}$) le corresponde el $A$-homomorfismo $\Gamma(h):M\to\Gamma(S,u_*(\sh{F}))=\Gamma(X,\sh{F})$.
\end{proof}

\begin{env}[1.8.8]
\label{ega2:addendum:I.1.8.8-fr}
Con la notación de \sref{ega2:addendum:I.1.8.4-fr}, se deduce claramente de
(\textbf 0, 4.1.1) que todo morfismo de espacios anillados
$(\psi,\theta):X\to S'$ equivale al dato de un homomorfismo de
anillos $A\to\Gamma(X,\sh{O}_X)$. Por tanto, la
Proposición~\sref{ega2:addendum:I.1.8.1-fr} puede interpretarse como una biyección
canónica
$\Hom(X,S)\xrightarrow{\sim}\Hom(X,S')$
donde en el segundo miembro se entienden morfismos de espacios
anillados, pues $A$ no es en general un anillo local.

Más generalmente, sean $X,Y$ espacios localmente anillados, y sea
$(Y',A')$ el espacio anillado cuyo espacio subyacente consta de un
punto y cuyo haz de anillos $A'$ es el haz simple definido por
$\Gamma(Y,O_Y)$. Entonces \sref{ega2:addendum:I.1.8.2.1-fr} puede interpretarse
como una aplicación
\begin{equation}
\label{ega2:addendum:I.1.8.8.1-fr}
  \rho:\Hom_{\mathit{an}}(X,Y)\to\Hom(X,Y')
\tag{1.8.8.1}
\end{equation}
El Corolario~\sref{ega2:addendum:I.1.8.3-fr} dice que los esquemas afines se
caracterizan, entre los espacios localmente anillados, como aquellos
para los cuales la restricción de $\rho$ a $\Hom(X,Y)$,
\begin{equation}
\label{ega2:addendum:I.1.8.8.2-fr}
  \rho':\Hom(X,Y)\to\Hom(X,Y')
\tag{1.8.8.2}
\end{equation}
es \emph{biyectiva} para \emph{todo} espacio localmente anillado $X$. En un capítulo
posterior generalizaremos esta definición, asociando así a un espacio
anillado \emph{arbitrario} $Z$ (y no solamente a uno cuyo espacio subyacente
consta de un punto) un espacio localmente anillado que volveremos a
denotar por $\Spec(Z)$. Este será el punto de partida de una teoría
«relativa» de los preesquemas sobre espacios anillados arbitrarios,
que extiende los resultados del Capítulo~I.
\end{env}

\begin{env}[1.8.9]
\label{ega2:addendum:I.1.8.9-fr}
Se pueden considerar los pares $(X,\sh{F})$ formados por un espacio localmente anillado $X$ y un $\sh{O}_X$-módulo $\sh{F}$ como formando una categoría, siendo un \emph{morfismo} $(X,\sh{F})\to(Y,\sh{G})$ de esta categoría un par $(u,h)$ formado por un morfismo de
\oldpage[II]{221}
espacios localmente anillados $u:X\to Y$ y de un $u$-morfismo
$h:\sh{G}\to\sh{F}$ de módulos; estos
morfismos (para $(X,\sh{F})$ e $(Y,\sh{G})$ fijados) forman un conjunto
que denotaremos por $\Hom((X,\sh{F}),(Y,\sh{G}))$; la aplicación
$(u,h)\to(\rho'(u),\Gamma(h))$ es una aplicación canónica
\begin{equation}
\label{ega2:addendum:I.1.8.9.1-fr}
\begin{split}
  \Hom((X,\sh{F}),(Y,\sh{G}))\to
  \Hom\bigl(
    &(\Gamma(Y,\sh{O}_Y),\Gamma(Y,\sh{G})),\\
    &(\Gamma(X,\sh{O}_X),\Gamma(X,\sh{F}))
  \bigr)
\end{split}
\tag{1.8.9.1}
\end{equation}
\emph{funtorial} en $(X,\sh{F})$ y $(Y,\sh{G})$, siendo el segundo miembro
el conjunto de los dihomomorfismos correspondientes a los anillos y a
los módulos considerados (\textbf 0, 1.0.2).
\end{env}

\begin{corollary}[1.8.10]
\label{ega2:addendum:I.1.8.10-fr}
Sean $Y$ un espacio localmente anillado y $\sh{G}$ un
$\sh{O}_Y$-módulo. Para que $Y$ sea un esquema afín y $\sh{G}$ un
$\sh{O}_Y$-módulo cuasicoherente, es necesario y suficiente que, para
toda pareja $(X,\sh{F})$ formada por un espacio localmente anillado $X$
y un $\sh{O}_X$-módulo $\sh{F}$, la aplicación canónica \textup{(1.8.8.1)} sea
biyectiva.
\end{corollary}

Dejamos al lector la tarea de desarrollar el razonamiento, modelado
sobre el de \sref{ega2:addendum:I.1.8.3-fr}, en el cual se utilizan \sref{ega2:addendum:I.1.8.1-fr} y
\sref{ega2:addendum:I.1.8.7-fr}.

\begin{remark}[1.8.11]
\label{ega2:addendum:I.1.8.11-fr}
Los enunciados (1.7.3), (1.7.4) y (2.2.4), así como la
definición (1.6.1), son casos particulares de \sref{ega2:addendum:I.1.8.1-fr}; del
mismo modo, (2.2.5) se deduce de \sref{ega2:addendum:I.1.8.7-fr}. La Proposición
(1.8.7) proporciona también (1.6.3) (y, por tanto, (1.6.4)) como
caso particular. En efecto, si $X$ es afín y $\Gamma(X,\sh{F})=N$, los
funtores
$M\to
  \Hom_{\sh{O}_S}(\widetilde M,u_*(\widetilde N))$ y $M\to
  \Hom_{\sh{O}_S}(\widetilde M,(N_{[\vphi]})\supertilde)$
(donde $\vphi:A\to\Gamma(X,\sh{O}_X)$ corresponde a $u$) son
isomorfos por \sref{ega2:addendum:I.1.8.7-fr} y (1.3.8). Finalmente, (1.6.5) (y,
por tanto, (1.6.6)) se deduce de \sref{ega2:addendum:I.1.8.6-fr} y del hecho de que,
para todo $f\in A$, los $A_f$-módulos
$N'\otimes_{A'}A_f$ y $(N'\otimes_{A'}A)_f$
(con las notaciones de (1.6.5)) son canónicamente isomorfos.
\end{remark}

\paragraph{Insértese después de (I, 3.2.8):}

\begin{env}[3.2.9]
\label{ega2:addendum:I.3.2.9-fr}
De \sref{ega2:addendum:I.1.8.1-fr} se deduce que (3.2.2) puede precisarse del modo
siguiente. El preesquema
$Z=\Spec(B\otimes_A C)$
no solo es un producto de $X=\Spec(B)$ e $Y=\Spec(C)$ en la categoría
de los $S$-\emph{preesquemas}, sino también en la categoría de los \emph{espacios
localmente anillados} \emph{sobre} $S$ (con la definición de los
$S$-morfismos modelada sobre (2.5.2)). La demostración de (3.2.6)
muestra asimismo que, para cualesquiera $S$-preesquemas $X,Y$, el
preesquema $X\times_S Y$ no solo es el producto de $X$ y $Y$ en la categoría de los
$S$-preesquemas, sino también en la categoría de los espacios
localmente anillados sobre el preesquema $S$.
\end{env}

\paragraph{(I, 4.4.5)}
En la línea 14 contando desde abajo de la p.~126, sustitúyase $B$ por
$A$; en la línea 13 contando desde abajo, sustitúyase «álgebra sobre
$A$» por «álgebra sobre $B$».

\paragraph{(I, 5.3.5)}
En la línea 2 contando desde abajo de la p.~132, sustitúyase $Y\to S$
por $g:Y\to S$.

\paragraph{(I, 6.3.2.1)}
En la línea 6 contando desde abajo de la p.~144, léase:
$D(g_i)\subset W$.

\paragraph{(I, 7.3.8)}
Sustitúyase el texto actual, que es erróneo, por el siguiente:

Sean $X,Y$ dos preesquemas \emph{íntegros}; entonces $\sh{R}(X)$ y
$\sh{R}(Y)$ son, respectivamente, un $\sh{O}_X$-módulo
cuasicoherente y un $\sh{O}_Y$-módulo cuasicoherente (7.3.3). Si
$f:X\to Y$ es un morfismo \emph{dominante}, existe un homomorfismo canónico
de $\sh{O}_X$-módulos
\oldpage[II]{222}
\begin{equation}
\label{ega2:addendum:I.7.3.8.1-fr}
  \tau:f^*(\sh{R}(Y))\to\sh{R}(X).
\tag{7.3.8.1}
\end{equation}

Supongamos primero que $X=\Spec(A)$ e $Y=\Spec(B)$ son afines, con
$A$ y $B$ anillos íntegros. Entonces $f$ corresponde a un homomorfismo
inyectivo $B\to A$, que se prolonga en un monomorfismo $L\to K$ del
cuerpo de fracciones $L$ de $B$ en el cuerpo de fracciones $K$ de $A$.
El homomorfismo \sref{ega2:addendum:I.7.3.8.1-fr} corresponde entonces al homomorfismo
canónico
$L\otimes_B A\to K$
de (1.6.5).

En el caso general, para toda pareja de abiertos afines no vacíos
$U\subset X$, $V\subset Y$ tales que $f(U)\subset V$, la construcción
anterior define un homomorfismo $\tau_{U,V}$. Si $U'\subset U$,
$V'\subset V$ y $f(U')\subset V'$, entonces $\tau_{U,V}$ prolonga a
$\tau_{U',V'}$; la afirmación se obtiene por pegado. Si $x,y$ son los
puntos genéricos de $X$ y $Y$, respectivamente, entonces $f(x)=y$ y
\[
  \bigl(f^*(\sh{R}(Y))\bigr)_x
  =\sh{O}_y\otimes_{\sh{O}_y}\sh{O}_x
  =\sh{O}_x
\]
por (\textbf 0, 4.3.1) de modo que $\tau_x$ es un \emph{isomorfismo}.

\paragraph{(I, 9.5.2)}
En las tres últimas líneas de la p.~176 y las cinco primeras de la
p.~177, sustitúyase en todas partes $X'$ por $Y'$ y $X''$ por $Y''$.

\paragraph{(I, 10.5.4)}
En la línea 11 contando desde abajo de la p.~187, sustitúyase (10.3.4)
por (10.3.5).

\paragraph{(I, 10.6.3)}
En la línea 3 contando desde abajo de la p.~189, sustitúyase $X$ por
$\mathfrak X$.

\paragraph{(I, 10.14.2)}
En la línea 5 de la p.~210, sustitúyase «$\sh{O}_{\mathfrak X}$-\emph{módulo
coherente}» por «\emph{ideal coherente de} $\sh{O}_{\mathfrak X}$»; en la
línea 7, sustitúyase «$\sh{O}_{\mathfrak X}$-módulos coherentes» por
«ideales coherentes de $\sh{O}_{\mathfrak X}$».
